\documentclass[10pt]{extarticle}
\usepackage[a4paper,margin=1.5cm]{geometry}
\usepackage[T1]{fontenc}
\usepackage{amsmath,amssymb}
\usepackage{graphicx}
\usepackage[export]{adjustbox}
\usepackage{multicol}
\usepackage{enumitem}
\usepackage{booktabs,array,tabularx}
\setlength{\columnsep}{0.7cm}
\setlength{\parindent}{0pt}
\setlength{\parskip}{3pt}
\pagestyle{empty}
\makeatletter
% figure registry (figdefs.tex is generated by the extract step)
\newcommand{\qfigdef}[4]{\@namedef{qf@file@#1}{#2}\@namedef{qf@w@#1}{#3}\@namedef{qf@wide@#1}{#4}}
\newcommand{\qf@check}[1]{\@ifundefined{qf@file@#1}{\PackageError{score2}{Unknown figure id #1}{}}{}}
\newcommand{\qf@img}[4][]{\includegraphics[width=#2\dimexpr\@nameuse{qf@w@#4}pt\relax,max width=#3,max height=0.45\textheight,#1]{\@nameuse{qf@file@#4}}}
% Full-width questions: finalize wraps every question that contains a wide figure in \qfullbegin...\qfullend,
% so the WHOLE question (stem, figures, options) is set across the page in reading order, then the two columns
% resume. The question's head - its number, stem and first figure (\qfullsplit marks the end of the head) -
% is kept on one page, so a number is never left at a page foot with its figure overleaf; the rest may break
% across pages normally, so no big gaps are left. Without a \qfullsplit, the whole question is kept together
% if it is under 45% of a page. Never write these yourself.
\newif\ifqf@full
\newif\ifqf@boxed
\newsavebox{\qf@box}
\newcommand{\qf@open}{\global\qf@fulltrue\par\global\qf@boxedtrue
  \setbox\qf@box\vbox\bgroup\hsize\textwidth\linewidth\textwidth}
\newcommand{\qf@emit}[1]{\par\egroup\global\qf@boxedfalse
  \ifdim\dimexpr\ht\qf@box+\dp\qf@box\relax<#1\textheight\noindent\usebox{\qf@box}\par\else\unvbox\qf@box\fi}
\newcommand{\qfullsplit}{\ifqf@boxed\qf@emit{0.85}\fi}
\newcommand{\qf@close}{\ifqf@boxed\qf@emit{0.45}\fi\global\qf@fullfalse}
\newcommand{\qfullbegin}{\end{multicols}\qf@open}
\newcommand{\qfullend}{\qf@close\begin{multicols}{2}}
\newcommand{\qfullstart}{\qf@open}
\newcommand{\qfulldone}{\qf@close}
% \qfig[scale]{id}: block figure, centred, directly where it is written (size set by the tool).
% Fallback if a wide figure is not inside a full-width question: break out of the columns in place - no floats.
\newcommand{\qfig}[2][1]{\qf@check{#2}%
  \ifqf@full\qf@blk{#1}{#2}\else
  \ifnum\@nameuse{qf@wide@#2}=1 \expandafter\@firstoftwo\else\expandafter\@secondoftwo\fi
  {\end{multicols}\begin{center}\qf@img{#1}{\textwidth}{#2}\end{center}\begin{multicols}{2}}%
  {\qf@blk{#1}{#2}}\fi}
\newcommand{\qf@blk}[2]{\par\smallskip{\centering\qf@img{#1}{\linewidth}{#2}\par}\smallskip}
% \qfigi[scale]{id}: inline figure (inside option lists / tables / minipages); never breaks columns.
\newcommand{\qfigi}[2][1]{\qf@check{#2}\qf@img[valign=c]{#1}{\linewidth}{#2}}
\makeatother
% structure macros - lint reads these to check completeness and figure placement
\newcommand{\papersec}[1]{\par\bigskip{\centering\Large\bfseries PAPER-#1\par}\medskip}
\newcommand{\partsec}[2]{\par\medskip\noindent{\large\bfseries PART-#1: #2}\par\smallskip}
\newcommand{\sect}[1]{\par\smallskip\noindent\textbf{#1}\par}
\newcommand{\qn}[1]{\par\smallskip\noindent\textbf{#1.}~}
\newcommand{\paraq}[2]{\par\smallskip\noindent\textbf{Paragraph for Questions #1 and #2}\par}
\newenvironment{qoptions}{\begin{itemize}[leftmargin=1.8em,itemsep=1pt,topsep=2pt,label={}]}{\end{itemize}}
\newcommand{\opt}[2]{\item[(#1)] #2}
% figure registry (generated by score2_tools.py finalize)
\qfigdef{p1_phy_s1_q01_1}{Score Advanced - Test 32 (08 May 2024)-img/p1_phy_s1_q01_1.png}{114.8}{0}
\qfigdef{p1_phy_s1_q03_1}{Score Advanced - Test 32 (08 May 2024)-img/p1_phy_s1_q03_1.png}{87.7}{0}
\qfigdef{p1_phy_s1_q04_1}{Score Advanced - Test 32 (08 May 2024)-img/p1_phy_s1_q04_1.png}{94.5}{0}
\qfigdef{p1_phy_s1_q05_1}{Score Advanced - Test 32 (08 May 2024)-img/p1_phy_s1_q05_1.png}{136.9}{0}
\qfigdef{p1_phy_s2_q07_1}{Score Advanced - Test 32 (08 May 2024)-img/p1_phy_s2_q07_1.png}{122.9}{0}
\qfigdef{p1_phy_s2_q08_1}{Score Advanced - Test 32 (08 May 2024)-img/p1_phy_s2_q08_1.png}{59.2}{0}
\qfigdef{p1_phy_s2_q09_1}{Score Advanced - Test 32 (08 May 2024)-img/p1_phy_s2_q09_1.png}{78.9}{0}
\qfigdef{p1_phy_s2_q10_1}{Score Advanced - Test 32 (08 May 2024)-img/p1_phy_s2_q10_1.png}{245.2}{0}
\qfigdef{p1_phy_s2_q11_1}{Score Advanced - Test 32 (08 May 2024)-img/p1_phy_s2_q11_1.png}{81.3}{0}
\qfigdef{p1_phy_s2_q12_1}{Score Advanced - Test 32 (08 May 2024)-img/p1_phy_s2_q12_1.png}{139.6}{0}
\qfigdef{p1_phy_s3_q01_1}{Score Advanced - Test 32 (08 May 2024)-img/p1_phy_s3_q01_1.png}{91.6}{0}
\qfigdef{p1_phy_s3_q02_1}{Score Advanced - Test 32 (08 May 2024)-img/p1_phy_s3_q02_1.png}{119.5}{0}
\qfigdef{p1_phy_s3_q04_1}{Score Advanced - Test 32 (08 May 2024)-img/p1_phy_s3_q04_1.png}{117.8}{0}
\qfigdef{p1_phy_s3_q05_1}{Score Advanced - Test 32 (08 May 2024)-img/p1_phy_s3_q05_1.png}{115.3}{0}
\qfigdef{p1_phy_s3_q06_1}{Score Advanced - Test 32 (08 May 2024)-img/p1_phy_s3_q06_1.png}{110.2}{0}
\qfigdef{p1_chem_s1_q03_1}{Score Advanced - Test 32 (08 May 2024)-img/p1_chem_s1_q03_1.png}{81.8}{0}
\qfigdef{p1_chem_s1_q04_1}{Score Advanced - Test 32 (08 May 2024)-img/p1_chem_s1_q04_1.png}{245.2}{0}
\qfigdef{p1_chem_s1_q04_A}{Score Advanced - Test 32 (08 May 2024)-img/p1_chem_s1_q04_A.png}{68.8}{0}
\qfigdef{p1_chem_s1_q04_B}{Score Advanced - Test 32 (08 May 2024)-img/p1_chem_s1_q04_B.png}{52.4}{0}
\qfigdef{p1_chem_s1_q04_C}{Score Advanced - Test 32 (08 May 2024)-img/p1_chem_s1_q04_C.png}{47.0}{0}
\qfigdef{p1_chem_s1_q04_D}{Score Advanced - Test 32 (08 May 2024)-img/p1_chem_s1_q04_D.png}{43.9}{0}
\qfigdef{p1_chem_s2_q07_1}{Score Advanced - Test 32 (08 May 2024)-img/p1_chem_s2_q07_1.png}{80.4}{0}
\qfigdef{p1_chem_s2_q08_A}{Score Advanced - Test 32 (08 May 2024)-img/p1_chem_s2_q08_A.png}{101.1}{0}
\qfigdef{p1_chem_s2_q08_B}{Score Advanced - Test 32 (08 May 2024)-img/p1_chem_s2_q08_B.png}{99.6}{0}
\qfigdef{p1_chem_s2_q08_C}{Score Advanced - Test 32 (08 May 2024)-img/p1_chem_s2_q08_C.png}{106.6}{0}
\qfigdef{p1_chem_s2_q08_D}{Score Advanced - Test 32 (08 May 2024)-img/p1_chem_s2_q08_D.png}{102.3}{0}
\qfigdef{p1_chem_s2_q09_R}{Score Advanced - Test 32 (08 May 2024)-img/p1_chem_s2_q09_R.png}{304.2}{1}
\qfigdef{p1_chem_s2_q09_3}{Score Advanced - Test 32 (08 May 2024)-img/p1_chem_s2_q09_3.png}{52.6}{0}
\qfigdef{p1_chem_s2_q10_A}{Score Advanced - Test 32 (08 May 2024)-img/p1_chem_s2_q10_A.png}{25.5}{0}
\qfigdef{p1_chem_s2_q10_B}{Score Advanced - Test 32 (08 May 2024)-img/p1_chem_s2_q10_B.png}{31.3}{0}
\qfigdef{p1_chem_s2_q10_C}{Score Advanced - Test 32 (08 May 2024)-img/p1_chem_s2_q10_C.png}{46.6}{0}
\qfigdef{p1_chem_s2_q10_D}{Score Advanced - Test 32 (08 May 2024)-img/p1_chem_s2_q10_D.png}{115.4}{0}
\qfigdef{p1_chem_s3_q03_1}{Score Advanced - Test 32 (08 May 2024)-img/p1_chem_s3_q03_1.png}{157.5}{0}
\qfigdef{p1_chem_s3_q04_1}{Score Advanced - Test 32 (08 May 2024)-img/p1_chem_s3_q04_1.png}{114.2}{0}
\qfigdef{p1_chem_s3_q04_2}{Score Advanced - Test 32 (08 May 2024)-img/p1_chem_s3_q04_2.png}{148.4}{0}
\qfigdef{p1_chem_s3_q04_3}{Score Advanced - Test 32 (08 May 2024)-img/p1_chem_s3_q04_3.png}{121.0}{0}
\qfigdef{p1_chem_s3_q04_4}{Score Advanced - Test 32 (08 May 2024)-img/p1_chem_s3_q04_4.png}{159.4}{0}
\qfigdef{p1_chem_s3_q05_1}{Score Advanced - Test 32 (08 May 2024)-img/p1_chem_s3_q05_1.png}{218.0}{0}
\qfigdef{p2_phy_s1_q01_1}{Score Advanced - Test 32 (08 May 2024)-img/p2_phy_s1_q01_1.png}{119.3}{0}
\qfigdef{p2_phy_s1_q03_1}{Score Advanced - Test 32 (08 May 2024)-img/p2_phy_s1_q03_1.png}{126.1}{0}
\qfigdef{p2_phy_s1_q04_1}{Score Advanced - Test 32 (08 May 2024)-img/p2_phy_s1_q04_1.png}{53.2}{0}
\qfigdef{p2_phy_s2_q05_1}{Score Advanced - Test 32 (08 May 2024)-img/p2_phy_s2_q05_1.png}{103.5}{0}
\qfigdef{p2_phy_s2_q08_1}{Score Advanced - Test 32 (08 May 2024)-img/p2_phy_s2_q08_1.png}{90.8}{0}
\qfigdef{p2_phy_s2_q09_1}{Score Advanced - Test 32 (08 May 2024)-img/p2_phy_s2_q09_1.png}{190.8}{0}
\qfigdef{p2_phy_s2_q10_1}{Score Advanced - Test 32 (08 May 2024)-img/p2_phy_s2_q10_1.png}{158.9}{0}
\qfigdef{p2_phy_s3_para1_2_1}{Score Advanced - Test 32 (08 May 2024)-img/p2_phy_s3_para1_2_1.png}{65.6}{0}
\qfigdef{p2_phy_s4_q07_1}{Score Advanced - Test 32 (08 May 2024)-img/p2_phy_s4_q07_1.png}{99.0}{0}
\qfigdef{p2_phy_s4_q08_1}{Score Advanced - Test 32 (08 May 2024)-img/p2_phy_s4_q08_1.png}{114.1}{0}
\qfigdef{p2_phy_s4_q09_1}{Score Advanced - Test 32 (08 May 2024)-img/p2_phy_s4_q09_1.png}{128.8}{0}
\qfigdef{p2_chem_s1_q01_1}{Score Advanced - Test 32 (08 May 2024)-img/p2_chem_s1_q01_1.png}{107.2}{0}
\qfigdef{p2_chem_s1_q01_A}{Score Advanced - Test 32 (08 May 2024)-img/p2_chem_s1_q01_A.png}{49.7}{0}
\qfigdef{p2_chem_s1_q01_B}{Score Advanced - Test 32 (08 May 2024)-img/p2_chem_s1_q01_B.png}{51.1}{0}
\qfigdef{p2_chem_s1_q01_C}{Score Advanced - Test 32 (08 May 2024)-img/p2_chem_s1_q01_C.png}{62.7}{0}
\qfigdef{p2_chem_s1_q01_D}{Score Advanced - Test 32 (08 May 2024)-img/p2_chem_s1_q01_D.png}{50.8}{0}
\qfigdef{p2_chem_s2_q07_1}{Score Advanced - Test 32 (08 May 2024)-img/p2_chem_s2_q07_1.png}{215.8}{0}
\qfigdef{p2_chem_s2_q07_A}{Score Advanced - Test 32 (08 May 2024)-img/p2_chem_s2_q07_A.png}{95.5}{0}
\qfigdef{p2_chem_s2_q07_B}{Score Advanced - Test 32 (08 May 2024)-img/p2_chem_s2_q07_B.png}{98.9}{0}
\qfigdef{p2_chem_s2_q07_C}{Score Advanced - Test 32 (08 May 2024)-img/p2_chem_s2_q07_C.png}{98.9}{0}
\qfigdef{p2_chem_s2_q07_D}{Score Advanced - Test 32 (08 May 2024)-img/p2_chem_s2_q07_D.png}{78.2}{0}
\qfigdef{p2_chem_s2_q08_A}{Score Advanced - Test 32 (08 May 2024)-img/p2_chem_s2_q08_A.png}{83.7}{0}
\qfigdef{p2_chem_s2_q08_B}{Score Advanced - Test 32 (08 May 2024)-img/p2_chem_s2_q08_B.png}{101.9}{0}
\qfigdef{p2_chem_s2_q08_C}{Score Advanced - Test 32 (08 May 2024)-img/p2_chem_s2_q08_C.png}{127.6}{0}
\qfigdef{p2_chem_s2_q08_D}{Score Advanced - Test 32 (08 May 2024)-img/p2_chem_s2_q08_D.png}{119.0}{0}
\qfigdef{p2_chem_s3_para1_2_1}{Score Advanced - Test 32 (08 May 2024)-img/p2_chem_s3_para1_2_1.png}{309.4}{1}
\qfigdef{p2_chem_s4_q09_1}{Score Advanced - Test 32 (08 May 2024)-img/p2_chem_s4_q09_1.png}{166.5}{0}
\raggedcolumns
\makeatletter\@beginparpenalty=10000\makeatother
\newcolumntype{L}[1]{>{\raggedright\arraybackslash}p{#1}}
\begin{document}
\begin{center}{\Large\bfseries Score Advanced --- Test 32 (08 May 2024)}\end{center}
\begin{multicols}{2}
\raggedcolumns
\papersec{1}
\partsec{1}{PHYSICS}
\sect{SECTION-I(i)}
\qn{1} A taut string is made of two segments. To the left of point `A', it has linear mass density of $\mu$ kg/m and to the right of A, its linear mass density is 4$\mu$ kg/m. A sinusoidal wave pulse is travelling towards right.\par
\qfig{p1_phy_s1_q01_1}
Which of the following is correct ?
\begin{qoptions}
\opt{A}{Ratio of amplitude of reflected and transmitted wave is $\dfrac14$}
\opt{B}{Ratio of amplitude of reflected and transmitted wave is 2}
\opt{C}{Ratio of transmitted power and reflected power is 8}
\opt{D}{Ratio of transmitted power and reflected power is 16}
\end{qoptions}
\qn{2} One mole of a monoatomic gas is heated in such a way that its molar specific heat is 2R. During the heating, the volume of the gas is doubled. By what factor does the temperature of the gas change ?
\begin{qoptions}
\opt{A}{$\dfrac1{\sqrt2}$}
\opt{B}{$\sqrt2$}
\opt{C}{2}
\opt{D}{4}
\end{qoptions}
\qn{3} The figure shows a uniform slab of mass 7 kg. Find the moment of inertia about the axis normal to the plane of drawing and passing through the point O. Given OA = OC, BC = 10 cm and AB = 2cm.\par
\qfig{p1_phy_s1_q03_1}
\begin{qoptions}
\opt{A}{84 kg cm$^2$}
\opt{B}{168 kg cm$^2$}
\opt{C}{188 kg cm$^2$}
\opt{D}{198 kg cm$^2$}
\end{qoptions}
\qn{4} An observer is standing at a point O, at a distance of 100 cm from a convex lens of focal length 50 cm. A plane mirror is placed behind the lens at a distance of 150 cm from the lens. The mirror now starts moving towards right with a velocity of 10 cm/s. What will be the magnitude of velocity of her own image as seen by the observer, at the moment when the mirror just starts moving ?\par
\qfig{p1_phy_s1_q04_1}
\begin{qoptions}
\opt{A}{$\dfrac{200}9$ cm/s}
\opt{B}{$\dfrac{200}3$ cm/s}
\opt{C}{$\dfrac{20}9$ cm/s}
\opt{D}{$\dfrac{10}9$ cm/s}
\end{qoptions}
\qn{5} Charge distribution on a rod bent as a semicircular arc of radius R follows a function $\lambda$ = $\lambda_0\sin\theta$. Here $\lambda$ is line charge density at a point on the rod, $\lambda_0$ is a positive constant and $\theta$ is angular position of the point as shown in the figure. V$_0$, V$_\mathrm P$ and V$_\mathrm Q$ represents the electrostatic potential at points O (centre point), P and Q respectively as shown in figure. Which of the following is correct?\par
\qfig{p1_phy_s1_q05_1}
\begin{qoptions}
\opt{A}{(V$_\mathrm P$ = V$_\mathrm Q$) $<$ V$_0$}
\opt{B}{(V$_\mathrm P$ = V$_\mathrm Q$) $>$ V$_0$}
\opt{C}{V$_\mathrm P$ = V$_\mathrm Q$ = V$_0$}
\opt{D}{V$_\mathrm P$ $>$ V$_\mathrm Q$ $>$ V$_0$}
\end{qoptions}
\qn{6} Nuclei of radioactive element A are being produced at rate t$^2$ at time t. The element A has decay constant equal to $\lambda$. If N is the number of active nuclei of element A at any time t then $\dfrac{dN}{dt}$ was found to be minimum at t = t$_0$. Find N at time t = t$_0$.
\begin{qoptions}
\opt{A}{$\dfrac{2\mathrm t_0-2\lambda\mathrm t_0^2}{\lambda^2}$}
\opt{B}{$\dfrac{2\mathrm t_0-\lambda\mathrm t_0^2}{\lambda^2}$}
\opt{C}{$\dfrac{\mathrm t_0-\lambda\mathrm t_0^2}{\lambda^2}$}
\opt{D}{$\dfrac{2\mathrm t_0-\lambda\mathrm t_0^2}{2\lambda^2}$}
\end{qoptions}
\sect{SECTION-I(ii)}
\qn{7} A ring shaped tube has uniform cross sectional area and its entire volume is 2V$_0$. The tube is well insulated from the surrounding. Inside the tube there is an adiabatic fixed wall S$_1$ and another movable adiabatic partition S$_2$. Initially, the movable partition is diametrically opposite to S$_1$ and the two halves of the tube have equal amount of an ideal gas ($\gamma$ = 1.5) at same pressure P$_0$. Now, a heater H is switched on which supplies heat slowly to one of the chambers. Heater is kept on till the partition S$_2$ moves through the quarter of the circle. At this position the heater is switched off and the partition S$_2$ remains in equilibrium. Neglect any friction as well as heat loss to the surrounding through the walls of the tube. Which of the following options are correct ?\par
\qfig{p1_phy_s2_q07_1}
\begin{qoptions}
\opt{A}{Final pressure in chamber `2' is 2P$_0$}
\opt{B}{Final pressure in chamber `2' is 2$\sqrt2$P$_0$}
\opt{C}{Work done on the gas in chamber `2' by the gas in chamber `1' is (2$\sqrt2$ $-$ 1)P$_0$V$_0$}
\opt{D}{Heat supplied by the heater to the gas is 4P$_0$V$_0$(2$\sqrt2$ $-$ 1)}
\end{qoptions}
\qn{8} A uniformly charged non conducting rod is suspended vertically at its end. The rod can swing freely in the vertical plane without any friction. The linear charge density on the rod is $\lambda$ C/m and it has a uniformly distributed mass of M. Length of the rod is L. A uniform horizontal electric field (E) is switched on in the region. Select the correct option/options :-\par
\qfig{p1_phy_s2_q08_1}
\begin{qoptions}
\opt{A}{If E = $\dfrac{\mathrm{mg}}{\lambda\mathrm L}$, then rod just manages to make it horizontal}
\opt{B}{If E = $\dfrac{\mathrm{mg}}{2\lambda\mathrm L}$, then rod just manages to make it horizontal}
\opt{C}{If E = $\dfrac{\mathrm{Mg}}{\lambda\mathrm L}$, the rod has maximum kinetic energy at an angle 45$^\circ$ with vertical}
\opt{D}{If E = $\dfrac{\mathrm{Mg}}{2\lambda\mathrm L}$, the rod has maximum kinetic energy at an angle 45$^\circ$ with vertical}
\end{qoptions}
\qn{9} A small ball of mass m is released from rest from a height h above the liquid surface. The density of ball is $\dfrac\rho3$ : (assume no splash)\par
\qfig{p1_phy_s2_q09_1}
\begin{qoptions}
\opt{A}{If h $\ge$ H, then the ball will be able to enter the liquid of higher density}
\opt{B}{If h $\ge$ $\dfrac{7\mathrm H}2$, then the ball will be able to strike the bottom of the vessel}
\opt{C}{If h = 2H, then the ball will cross the interface of two liquids with kinetic energy = mgH}
\opt{D}{If h = 0 then ball will float on the surface.}
\end{qoptions}
\qn{10} Consider three fixed surfaces shown in the figure. Three blocks each of mass m are released from rest from top of the three surfaces. All blocks reach ground with same speed. Length of path travelled by the blocks is same for second and third surface. If coefficient of friction of three surfaces are $\mu_1$, $\mu_2$ and $\mu_3$ respectively then\par
\qfig{p1_phy_s2_q10_1}
\begin{qoptions}
\opt{A}{$\mu_1$ = $\mu_2$}
\opt{B}{$\mu_1$ $>$ $\mu_2$}
\opt{C}{$\mu_2$ = $\mu_3$}
\opt{D}{$\mu_2$ $<$ $\mu_3$}
\end{qoptions}
\qn{11} In the circuit shown in figure R and V$_0$ are constant, while r can change, taking up all values. What can be the equivalent resistance of the circuit ?\par
\qfig{p1_phy_s2_q11_1}
\begin{qoptions}
\opt{A}{1.6 R}
\opt{B}{1.4 R}
\opt{C}{1.8 R}
\opt{D}{2.2 R}
\end{qoptions}
\qn{12} There are two inclined planes, as shown in the figure. A ball strikes the inclined plane AB vertically with velocity v = $\sqrt{24\,\mathrm g}$ at height h = $\left(1+2\sqrt3\right)$ m. The collision is inelastic. After this the ball moves such that it collides with the inclined plane BC perpendicular to it. Select the option(s).\par
\qfig{p1_phy_s2_q12_1}
\begin{qoptions}
\opt{A}{The co-efficient of restitution between the inclined plane AB and ball is $\dfrac1{\sqrt3}$}
\opt{B}{The co-efficient of restitution between the inclined plane AB and the balls is $\dfrac1{3\sqrt3}$}
\opt{C}{The time interval between two collision is $\sqrt{\dfrac8{\mathrm g}}$ s}
\opt{D}{The velocity of collision with the inclined plane BC is $\sqrt{6\,\mathrm g}\left(1+\dfrac1{\sqrt3}\right)$ m/s}
\end{qoptions}
\sect{SECTION-II}
\qn{1} A non-viscous incompressible liquid of mass m is filled fully inside a thin uniform spherical shell of mass m and radius R performing pure rolling on a rough horizontal surface. There are two point A and B inside the liquid on the vertical diameter separated by 2R as shown in the figure. Pressure difference between B and A is (P$_\mathrm B-$P$_\mathrm A$). At the given instant velocity of centre of mass of given system is V$_0$ and kinetic energy of this system is K. If the value of $\dfrac{\mathrm K}{\mathrm P_\mathrm B-\mathrm P_\mathrm A}$ is equal to $\dfrac{\mathrm n\pi\mathrm V_0^2\mathrm R^2}{9\mathrm g}$. Find n.\par
\qfig{p1_phy_s3_q01_1}
\qn{2} A block of mass `m' is given velocity towards another block of mass 6kg kept at rest on a frictionless horizontal floor some distance apart from a wall as shown in the figure. If all collisions are elastic, find minimum value of mass `m' (in kg) so that blocks will collide only once.\par
\qfig{p1_phy_s3_q02_1}
\qn{3} In a hydrogen like atom an electron is orbiting in an orbit having quantum number n. Its frequency of revolution is found to be 13.2 $\times$ 10$^{15}$ Hz. Energy required to free the electron from the atom from the above orbit is 54.4 eV. In time 7 nano-seconds the electron jumps back to orbit having quantum number $\dfrac{\mathrm n}2$.\par
If $\tau$ be the average torque acted on the electron during the above process, then find $\dfrac\tau3\times10^{27}$ in N$-$m. (Given $\dfrac h\pi$ = 2.1 $\times$ 10$^{-34}$ J$-$s).\par
Frequency of revolution of electron in the ground state of H-atom f$_0$ = 6.6 $\times$ 10$^{15}$ Hz and ionisation energy of H-bond E$_0$ = 13.6 eV
\qn{4} As shown in the figure a cylinder of mass M and radius R moves with constant speed v through a region of space that contains dust particles of mass m which are at rest. Number of particles in unit volume is n. The cylinder moves in a direction perpendicular to its axis. Assume that the particles do not interact with each other and that m $<<$ M. All collisions are elastic and surface of the cylinder is smooth. The force per unit length of the cylinder required to maintain the speed v constant is $\dfrac{\mathrm K\,\mathrm{nmRv}^2}3$. Find the value of K.\par
\qfig{p1_phy_s3_q04_1}
\qn{5} A positively charged particle is given initial velocity v$_0$ in XY plane as shown in diagram. Initially the particle was at origin and the region contains uniform, mutually perpendicular electric and magnetic field. Let $\vec{\mathrm E}$ = $-$E$\hat{\mathrm j}$ and $\vec{\mathrm B}$ = B$\hat{\mathrm k}$. The minimum magnitude of the magnetic field so that the charge particle crosses the origin again during its subsequent motion is B$_{\min}$ = $\dfrac{\mathrm n\pi\mathrm E}{\mathrm v_0}$. Find n\par
\qfig{p1_phy_s3_q05_1}
\qn{6} In the arrangement shown, two springs are attached to a light rod as shown. Length of O to A is $\dfrac{2\ell}3$. System is at equilibrium. The time period of oscillation of the block, if it is displaced slightly is found to be T = 2$\pi\sqrt{\dfrac{\alpha\mathrm m}{8\mathrm K}}$. Find the value of $\alpha$.\par
\qfig{p1_phy_s3_q06_1}
\partsec{2}{CHEMISTRY}
\sect{SECTION-I(i)}
\qn{1} A metal exists in two allotropic forms, $\alpha$ and $\beta$, having a face centred cubic and hexagonal closest packed lattice respectively. Identify the correct statement.
\begin{qoptions}
\opt{A}{Coordination number is higher for the $\alpha$-form.}
\opt{B}{Density is higher for the $\beta$-form.}
\opt{C}{Distance between closest octahedral voids is same for both the forms.}
\opt{D}{Unit cell volume is same for both the forms.}
\end{qoptions}
\qn{2} Consider two Van der waal gases A and B. For gas A, $\displaystyle\lim_{P\to0}\left(\dfrac{\partial Z}{\partial P}\right)_{T=300\,\mathrm K}>0$ and for gas B, $\displaystyle\lim_{P\to0}\left(\dfrac{\partial Z}{\partial P}\right)_{T=300\,\mathrm K}<0$. Identify the incorrect statement.
\begin{qoptions}
\opt{A}{Boyle temperature of gas A $<$ 300 K}
\opt{B}{Boyle temperature of gas B $>$ 300K.}
\opt{C}{Boyle temperature of both gases is 300 K}
\opt{D}{Critical temperature of B is higher than that of A}
\end{qoptions}
\qn{3} In which of the following cases, the most stable conformer is achiral?\par
\qfig{p1_chem_s1_q03_1}
\begin{qoptions}
\opt{A}{Only IV}
\opt{B}{I and IV}
\opt{C}{I, II and III}
\opt{D}{In all the cases}
\end{qoptions}
\qn{4}\par
\qfig{p1_chem_s1_q04_1}
product `D' is:
\begin{qoptions}
\opt{A}{\qfigi{p1_chem_s1_q04_A}}
\opt{B}{\qfigi{p1_chem_s1_q04_B}}
\opt{C}{\qfigi{p1_chem_s1_q04_C}}
\opt{D}{\qfigi{p1_chem_s1_q04_D}}
\end{qoptions}
\qn{5} Select the incorrect statement.
\begin{qoptions}
\opt{A}{The less stable tautomeric form of H$_3$PO$_3$ has three active hydrogens.}
\opt{B}{Both H$_3$PO$_3$ and H$_3$PO$_2$ disproportionate on heating to produce PH$_3$ and H$_3$PO$_4$.}
\opt{C}{In P$_4$O$_{10}$ all the O-P-O bond angles are identical whereas there are two different P-O bond lengths.}
\opt{D}{H$_3$PO$_3$ and H$_3$AsO$_3$ have different basicity.}
\end{qoptions}
\qn{6} Which of the following statement is correct regarding extraction of aluminium from bauxite ore:
\begin{qoptions}
\opt{A}{In Hall's process CO$_2$ is formed at the anode.}
\opt{B}{Baeyer's process is employed when the bauxite ore is rich in silica.}
\opt{C}{In Serpeck's process N$_2$ is produced due to hydrolysis of aluminium nitride.}
\opt{D}{Bauxite ore undergoes roasting before the leaching process.}
\end{qoptions}
\sect{SECTION-I(ii)}
\qn{7} 1 mole of He $\left(\mathrm C_{v,m}=\dfrac{3R}2\right)$ at 1 bar is compressed adiabatically against a constant external pressure, P$_\mathrm{ext}$ from an initial volume of 20 L. Assuming He to show ideal gas behaviour, identify the correct statement(s) for the process. (Take R = 2 cal mol$^{-1}$K$^{-1}$, ln3 = 1.1, ln5 = 1.6)\par
\qfig{p1_chem_s2_q07_1}
\begin{qoptions}
\opt{A}{The lowest possible final volume of the gas after this process is 8L.}
\opt{B}{If P$_\mathrm{ext}$ = 3 atm, then $\Delta$S = 0.8 calK$^{-1}$}
\opt{C}{$\Delta$S$_\mathrm{sys}$ is +ve and $\Delta$S$_\mathrm{Surroundings}$ = 0 for all values of P$_\mathrm{ext}$}
\opt{D}{If T$_1$, V$_1$, are the initial temperature and volume of the system, and T$_2$, V$_2$ are the final temperature and volume respectively, then $\ln\left(\dfrac{T_2}{T_1}\right)=\dfrac R{C_v}\ln\left(\dfrac{V_1}{V_2}\right)$}
\end{qoptions}
\qn{8} An azeotrop of HNO$_3$ and H$_2$O is formed with 38 mole percent of HNO$_3$. Identify the correct plot(s).\par
y $\to$ vapour phase mole fraction.\par
$\chi$ $\to$ liquid phase mole fraction.
\begin{qoptions}
\opt{A}{\qfigi{p1_chem_s2_q08_A}}
\opt{B}{\qfigi{p1_chem_s2_q08_B}}
\opt{C}{\qfigi{p1_chem_s2_q08_C}}
\opt{D}{\qfigi{p1_chem_s2_q08_D}}
\end{qoptions}
\qfullbegin
\qn{9}\par
\qfig{p1_chem_s2_q09_R}
\qfullsplit
Identify the correct statement(s):
\begin{qoptions}
\opt{A}{Structure of `E' is $\mathrm{CH_3-CH=CH-\overset{O}{\overset{\|}{C}}-H}$.}
\opt{B}{Structure of `B' and `C' can be differentiated by Fehling solution.}
\opt{C}{Structure of `E' is \qfigi{p1_chem_s2_q09_3}}
\opt{D}{Structure of `A' is Ph $-$ CH $=$ CH $-$ CH$_3$}
\end{qoptions}
\qfullend
\qn{10} Which of the following is/are not correctly matched?
\begin{qoptions}
\opt{A}{\qfigi{p1_chem_s2_q10_A} $\xrightarrow{\text{DCl}}$ Number of stereoisomers of intermediate and products are identical.}
\opt{B}{\qfigi{p1_chem_s2_q10_B} $\xrightarrow[\text{CCl}_4]{\text{Br}_2}$ Neither product nor intermediate exhibits optical activity.}
\opt{C}{\qfigi{p1_chem_s2_q10_C} $\xrightarrow[\text{Pd}]{\text{H}_2}$ The chair form of the product gives non-superimposable structure after ring flipping.}
\opt{D}{\qfigi{p1_chem_s2_q10_D} Product mixture upon fractional distillation gives 4 major fractions, out of which 3 are optically active.}
\end{qoptions}
\qn{11} Select the correct statement(s):
\begin{qoptions}
\opt{A}{[Cr(NH$_3$)$_5$SO$_4$]Cl will produce a white precipitate with AgNO$_3$ solution but not with BaSO$_4$ solution.}
\opt{B}{Octahedral complexes formed by coordinating strong field ligands with Fe$^{+3}$ will be diamagnetic in nature.}
\opt{C}{[Ni(en)$_3$]Cl$_2$, will show optical as well as geometrical isomerism.}
\opt{D}{K[PtCl$_3$C$_2$H$_4$] is an organometallic compound}
\end{qoptions}
\qn{12} A (is a compound of Ba) and B (is a sulphate) are white crystalline substances, highly soluble in water and decompose at higher temperatures. If an aqueous solution of A (slightly basic) is added to an aqueous solution of B (slightly acidic) a white precipitate C forms. The filtrate is neutral and gives a brown colour solution with an excess of acidified KI solution. When boiled, the filtrate evaporates without the appearance of any residue. White solid D can be prepared by the heating of A in the absence of air. The reaction of D with water gives a colourless solution. This solution, if kept in an open container, slowly precipitates a white solid E and leaves water. Upon prolonged exposure to moist air at room temperature, solid D is transformed into E as well. E can be converted back into D but decomposition above 1400$^\circ$C is required for this purpose. The reaction of B and D in water forms the precipitate C and is accompanied by a characteristic odour due to formation of a gas. A gives white precipitate with AgNO$_3$. Select the correct option(s).
\begin{qoptions}
\opt{A}{Compound A is on heating gives neutral, acidic and basic oxide}
\opt{B}{Compound B is (NH$_4$)$_2$SO$_4$.}
\opt{C}{Compound E is BaCO$_3$.}
\opt{D}{I$_2$ will precipitate as brown coloured.}
\end{qoptions}
\sect{SECTION-II}
\qn{1} Hydrogen peroxide in aqueous solution decomposes by a first order reaction to produce water and oxygen. The rate constant for this decomposition is 7.50 $\times$ 10$^{-4}$ s$^{-1}$. What quantity of heat (in J) is initially liberated per second from 2.0 L of solution that is 0.02M H$_2$O$_2$? Given $\Delta_f$H$^\circ$ of H$_2$O$_2$($l$) = $-$187.0 kJ/mol and $\Delta_f$H$^\circ$ H$_2$O ($l$) = $-$287.0 kJ/mol.
\qn{2} Calculate the pH of solution formed by addition of 1 mole each of HCl, CH$_3$COONa and HCOONa in enough water to make 1 L solution. The pK$_\mathrm a$ of HCOOH and CH$_3$COOH are 4 and 5, respectively.
\qn{3} In the given reactions:-\par
\qfig{p1_chem_s3_q03_1}
Number of correct statement(s) among the following are :-\par
(1) X and Y are enantiomers\par
(2) M and N are enantiomers\par
(3) A and B are diastereomers\par
(4) (X and Y) are identical with (A and B)\par
(5) (P and Q) are identical with (M and N)\par
(6) M and N are geometrical isomers
\qn{4} In how many of the below options, the first structure is more stable than the second.\par
\qfig{p1_chem_s3_q04_1}
\qfig{p1_chem_s3_q04_2}
\qfig{p1_chem_s3_q04_3}
\qfig{p1_chem_s3_q04_4}
\qn{5} If the reaction sequence given below is carried out with 4 mol of BrO$_3^-$ and just required amount of Br$^-$ and H$^+$, the amount of the product D formed (in g) is\_\_\_\_\_\_\_\_\_\_\_\_\_.\par
The yields of A, B, C, D are given in the parentheses.\par
\qfig{p1_chem_s3_q05_1}
\qn{6} The bond angle H-C-H in the formaldehyde molecule is approximately 120$^\circ$; the dipole moments for C-H and C = O bonds are $\mu_\mathrm{C-H}$ = 0.4 D and $\mu_\mathrm{C=O}$ = 2.3 D, respectively. The net dipole moment of the formaldehyde is X $\times$ 10$^{-\mathrm Y}$ Cm (in scientific notation). The value of X is\_\_\_\_\_\_\_\_\_.
\partsec{3}{MATHEMATICS}
\sect{SECTION-I(i)}
\qn{1} If $\alpha$ and $\beta$ are exactly two distinct real roots of the equation $(x^2+5x-b^2)(2x^2-ax-6)=0$ where a, b $\in$ R, (Given that equation has only real roots) then
\begin{qoptions}
\opt{A}{$\alpha^2+\beta^2=30$}
\opt{B}{$a^2+b^2=105$}
\opt{C}{$\alpha^2+\beta^2=19$}
\opt{D}{$a^2+b^2+\alpha^2+\beta^2=134$}
\end{qoptions}
\qn{2} If f(x) = x$^4$ $-$ 3x$^2$ + 7x $-$ 1 and g(x) = cot$^{-1}$(2x$^2$) + $\sin^{-1}\left(\dfrac1{1+x^2}\right)$ are two functions, then number of integers in the range of gof(x) is
\begin{qoptions}
\opt{A}{1}
\opt{B}{2}
\opt{C}{3}
\opt{D}{4}
\end{qoptions}
\qn{3} $\displaystyle\lim_{x\to0^+}\dfrac{(1+\tan x)^{2\cot x}-e^{2\cos x}}{1-\cos(\ln(1+\sin\sqrt x))}$
\begin{qoptions}
\opt{A}{$2e^2$}
\opt{B}{$2e$}
\opt{C}{$-2e$}
\opt{D}{$-2e^2$}
\end{qoptions}
\qn{4} $\displaystyle\int e^x\left[\dfrac{2\sin x}{(\sin x+\cos x)^2}-\dfrac{2\cos x}{(\sin x-\cos x)^2}\right]dx$ where `C' is the constant of integration.
\begin{qoptions}
\opt{A}{$e^x\left[\dfrac{2\sin x}{\sin^2x-\cos^2x}\right]+C$}
\opt{B}{$e^x\left[\dfrac{2\cos x}{\sin^2x-\cos^2x}\right]+C$}
\opt{C}{$e^x\left(\dfrac{\sin^2x-\cos^2x}{2\sin x}\right)+C$}
\opt{D}{$e^x\left(\dfrac{\sin^2x-\cos^2x}{2\cos x}\right)+C$}
\end{qoptions}
\qn{5} The value of S = $(\sin^3\theta+\cos^3\theta)+(\sin^33\theta+\cos^33\theta)+(\sin^35\theta+\cos^35\theta)+\dots\dots\dots\dots\dots+(\sin^319\theta+\cos^319\theta)$ when $\theta=\dfrac\pi{12}$ is $\dfrac{-a}{\sqrt b}$ where a \& b are realtively prime numbers, then $|\,\mathrm b-6\mathrm a\,|$ equals
\begin{qoptions}
\opt{A}{1}
\opt{B}{2}
\opt{C}{3}
\opt{D}{4}
\end{qoptions}
\qn{6} Consider f(x) = x$^3$ + ax$^2$ + bx +c where a, b, c $\in$ \{1, 2, 3, 4, 5, 6\}. The probability that f(x) = 0 has all real roots given that f(x) is non-decreasing function, is
\begin{qoptions}
\opt{A}{$\dfrac1{216}$}
\opt{B}{$\dfrac1{96}$}
\opt{C}{$\dfrac1{108}$}
\opt{D}{$\dfrac1{36}$}
\end{qoptions}
\sect{SECTION-I(ii)}
\qn{7} The vertices of triangle ABC are A(2, 0, 2); B($-$1, 1, 1) and C(1, $-$2, 4). The point D and E divide AB and CA in the ratio 1 : 2 respectively. Another point F is taken in space such that perpendicular drawn from F to the plane containing $\Delta$ABC, meets the plane at the point of intersection (point P) of the line segment CD \& BE. If the distance of F from the plane of triangle ABC is $\sqrt2$ units then
\begin{qoptions}
\opt{A}{$\dfrac{\text{area of }\Delta\mathrm{ABC}}{\text{area of }\Delta\mathrm{PBC}}=\dfrac{7\sqrt3}6$}
\opt{B}{Volume of tetrahedral ABCF $=\dfrac73$ cubic units}
\opt{C}{Equation of line AF is $\vec r=2(\hat i+\hat k)+\lambda(\hat i+2\hat j)$}
\opt{D}{Equation of line AF is $\vec r=2(\hat i+\hat k)+\lambda(\hat i-2\hat k)$}
\end{qoptions}
\qn{8} If chords of the ellipse $\dfrac{x^2}{a^2}+y^2=2$, passing through (a, $-$1) are bisected by the line x + y = 1, then
\begin{qoptions}
\opt{A}{least positive integral value of a is 2.}
\opt{B}{least positive integral value of a is 1.}
\opt{C}{sum of integral values of a in ($-$10, 5) is $-$8.}
\opt{D}{sum of negative integral values of a in ($-$20, $-$9) is $-$145.}
\end{qoptions}
\qn{9} Let f : R $\longrightarrow$ R be a differentiable function satisfying f(x + y) = f(x) + f(y) + (x + y) e$^{x+y}$ $-$ xe$^x$ $-$ ye$^y$ + 4xy $\forall$ x, y $\in$ R and f $'$(0) = 1, then which of the following option(s) is/are correct?
\begin{qoptions}
\opt{A}{f(x) = e$^{-x+3}$; x $>$ 0 has exactly one solution}
\opt{B}{f(2) $-$ f $'$(2) $<$ 0}
\opt{C}{3f(x) = sgn($-$x$^2$ + 3x $-$ 6) has exactly one solution in x $\in$ [2, 5]}
\opt{D}{f(3) $-$ 3e$^3$ = x$^2$ $-$ 6x + 5 has exactly two solution}
\end{qoptions}
\qn{10} $\displaystyle\int_0^{\pi}[\,\{x\}+[\cot x]\,]\,dx$ is greater than or equal to\par
[Note: \{z\} \& [z] denotes fractional part function and greatest integer function of z.]
\begin{qoptions}
\opt{A}{0}
\opt{B}{1}
\opt{C}{$-\dfrac\pi2$}
\opt{D}{$-\pi$}
\end{qoptions}
\qn{11} Let OABC is a rhombus where O is origin such that vertices A (z$_1$) and C (z$_3$) lie on the circle $|z|=2$ and the vertex B (z$_2$) lie on the circle $|z|=3$. Let Arg (z$_2-$ z$_1$) = $\dfrac\pi3$, then
\begin{qoptions}
\opt{A}{$z_1=\left(\dfrac{3-\sqrt7i}{3+\sqrt7i}\right)(1+\sqrt3i)$}
\opt{B}{$z_2=\left(\dfrac{6+6\sqrt3i}{3-\sqrt7i}\right)$}
\opt{C}{$z_3=(1+\sqrt3i)$}
\opt{D}{$z_1=\left(\dfrac{3-\sqrt7i}{3+\sqrt7i}\right)(1-\sqrt3i)$}
\end{qoptions}
\qn{12} Let f(x) = $-$ x$^3$ + px$^2$ + qx + 6sgn (x$^2$ + 2x + 5) where p, q $\in$ R \& sgn denotes signum function. If the largest possible interval in which f $'$(x) is positive is ($-$ 2, 1), then (p + q) is less than or equal to
\begin{qoptions}
\opt{A}{4}
\opt{B}{5}
\opt{C}{6}
\opt{D}{9}
\end{qoptions}
\sect{SECTION-II}
\qn{1} Let a plane P containing a line $\vec r=\left(\hat i+\hat j+\hat k\right)+\lambda\left(-\hat i+3\hat j+2\hat k\right)$ and perpendicular to 2x + 3y $-$ z = 2. If the area of projection triangle formed on plane P whose vertices are A$_r$ = $2r\hat i+3(-1)^r\hat j+r^2\hat k$ where r $\in$ N \& 1 $\leqslant$ r $\leqslant$ 3 is a$\sqrt b$ where a, b are relatively prime numbers, then $\left(\dfrac ba\right)$ is equal to
\qn{2} In a book with page numbers 1 to n, four consecutive pages are torn off and the sum of the numbers on the remaining pages is 1596, then find the sum of possible values of n. (4 consecutive pages means page numbers m, m+1, m+2 \& m+3, m$\in$ N )
\qn{3} From a fixed point A, three normals are drawn to the parabola y$^2$ = 4ax at the points P, Q and R. Two circles C$_1$ and C$_2$ be drawn on AP and AQ as diameter. If slope of the common chord of the circles C$_1$ and C$_2$ be m$_1$ and the slope of the tangent to the parabola at the point R be m$_2$, then m$_1$ $\times$ m$_2$ is equal to
\qn{4} Let S be the set containing all matrices A = [a$_{ij}$]$_{3\times3}$ having entries from the set \{0,1,2,4\} such that sum of all principal diagonal entries of AA$^\mathrm T$ = 17 then number of matrices in S is
\qn{5} If $\displaystyle\int_0^1\left(\left(6x^4-f(x)\right)\right)f(x)\,dx=1$, then the area of region enclosed by y = f (x), x = 0 and x = 2.
\qn{6} Let u = sin$^7$x, v = cos x. If $\left.\dfrac{d^2}{dx^2}\left(\dfrac{d^3u}{dv^3}\right)\right|_{x=\frac\pi8}=2^{\lambda_1}\cdot3^{\lambda_2}\cdot5^{\lambda_3}\cdot7^{\lambda_4}$, then the value of ($\lambda_1+\lambda_2+\lambda_3+\lambda_4$), is
\raggedcolumns
\papersec{2}
\partsec{1}{PHYSICS}
\sect{SECTION-I(i)}
\qn{1} A small block of mass m lies above a thin disk of total mass M, constant surface density $\sigma$, and radius R. A hole of radius R/2 is cut into the center of the disk, allowing the mass to pass though. The mass starts at a position along the central axis of the disk but is displaced a height $\Delta$h from the plane of the disk, where $\Delta$h $<<$ R. There are no other external forces at work on this system; i.e., the disk and block are out in deep space. Assume m $<<$ M. For the first order in $\Delta$h the block undergoes simple harmonic motion with respect to the centre of the plane of the disk, compute the period of oscillation of block.\par
\qfig{p2_phy_s1_q01_1}
\begin{qoptions}
\opt{A}{$T=\sqrt{\dfrac{2\pi R}{G\sigma}}$}
\opt{B}{$T=\sqrt{\dfrac{\pi R}{G\sigma}}$}
\opt{C}{$T=\sqrt{\dfrac{\pi R}{2G\sigma}}$}
\opt{D}{None of these}
\end{qoptions}
\qn{2} Consider a metal can placed coaxially inside a long solenoid. The metal can be modelled as a thin cylindrical shell of radius r. The magnetic field due to the solenoid varies with time t as kt where k is a positive constant. The resistivity of the can is $\rho$ and its breaking stress is $\sigma$. Ignore the self-inductance of the can and any currents induced in the flat surfaces of the can. Assume that the magnetic field produced by the induced currents is negligible. Compute the time at which the can breaks.
\begin{qoptions}
\opt{A}{$t=\dfrac{\sigma\rho}{r^2K^2}$}
\opt{B}{$t=\dfrac{2\sigma\rho}{r^2K^2}$}
\opt{C}{$t=\dfrac{\sigma\rho}{2r^2K^2}$}
\opt{D}{$t=\dfrac{\sigma\rho}{3r^2K^2}$}
\end{qoptions}
\qn{3} A motor is installed at the top of a pole rigidly fixed on a platform. A light rod of length r = 1 m is rigidly attached to the motor shaft at its one end and at the other end a small ball of mass m is attached. The rod can rotated in a vertical plane with the help of the motor. Total mass of the platform, the pole and the motor is $\eta$ = 4 times the mass of the ball. The motor rotates the rod at a constant angular velocity. The platform is placed on a horizontal surface, where coefficient of friction is $\mu$ = 1$\sqrt3$. At which minimum angular velocity $\omega_0$ of the rod, will the platform start sliding?\par
\qfig{p2_phy_s1_q03_1}
\begin{qoptions}
\opt{A}{$\omega_0$ = 2rad / s}
\opt{B}{$\omega_0$ = 3rad / s}
\opt{C}{$\omega_0$ = 5rad / s}
\opt{D}{None of these}
\end{qoptions}
\qn{4} A ring of mean radius r and cross-sectional area A is made of a perfectly conducting wire. Inductance L of the ring is so small that inertia of free electrons cannot be neglected in the current building process. The free electron density in the conductor is n, mass of an electron is m and modulus of charge on an electron is e. Initially the ring is placed in a uniform magnetic field with its plane parallel to induction vector $\vec{\mathrm B}$ of the field as shown in the figure. Find the current ring after it is turned through an angle 90$^\circ$.\par
\qfig{p2_phy_s1_q04_1}
\begin{qoptions}
\opt{A}{$i=\dfrac{B\pi r^2}{L+\frac{m\times\pi r}{ne^2A}}$}
\opt{B}{$i=\dfrac{B\pi r}{L+\frac{m\times\pi r}{ne^2A}}$}
\opt{C}{$i=\dfrac{B\pi r^2}{L+\frac{m\times2\pi r}{ne^2A}}$}
\opt{D}{$i=\dfrac{B\times2\pi r}{L+\frac{m\times\pi r^2}{ne^2A}}$}
\end{qoptions}
\sect{SECTION-I(ii)}
\qn{5}\par
\qfig{p2_phy_s2_q05_1}
For the given figure consider their mutual inductance is M then
\begin{qoptions}
\opt{A}{Leq $=\dfrac{L_1L_2-M^2}{L_1+L_2-2M}$ When mutual inductance is not opposing}
\opt{B}{Leq $=\dfrac{L_1L_2-M^2}{L_1+L_2+2M}$ When mutual inductance is not opposing.}
\opt{C}{Leq $=\dfrac{L_1L_2-M^2}{L_1+L_2+2M}$ When mutual inductance is opposing.}
\opt{D}{Leq $=\dfrac{L_1L_2-M^2}{2L_1+L_2+2M}$ When mutual inductance is opposing.}
\end{qoptions}
\qn{6} In a photoelectric effect experiment, a point source of light of power 40W emits mono-energetic photons of wavelength $\lambda_1$ that can just exit photoelectrons from an isolated metallic sphere of radius 1 cm, placed at a distance of 1 km from the light source. Now, three other sources of wavelength $\lambda_2$ ,$\lambda_3$ and $\lambda_4$ which are emitting the same number of photons as that of $\lambda_1$ brought near the source of $\lambda_1$. Assume photo efficiency of 10$^{-6}$ (Take $\lambda_1$ = 4960\AA, $\lambda_2$ = 4133.33\AA, $\lambda_3$ = 5000\AA, $\lambda_4$ = 7200\AA and hc = 12.400 eV.\AA)
\begin{qoptions}
\opt{A}{Number of photoelectrons emitted from the sphere per second is 2.5 $\times$ 10$^9$}
\opt{B}{The potential of the sphere when emission of photoelectron stops is 0.8 V.}
\opt{C}{The potential of the sphere when emission of photoelectrons stops is 0.5 V.}
\opt{D}{The time after which the emission of photoelectrons will stop is 1.39 $\times$10$^{-3}$ sec.}
\end{qoptions}
\qn{7} A dog running with a constant speed v is chasing a cat that is running with a constant velocity $\vec{\mathrm u}$. During the chase, the dog always heads towards the cat. At an instant, direction of motion of the dog makes angle $\theta$ with that of the cat and the distance between them is r.
\begin{qoptions}
\opt{A}{The magnitude of acceleration of dog at this instant is $\dfrac{uv\sin\theta}{r}$.}
\opt{B}{Radius of curvature of path followed by dog at this instant is $\dfrac{vr}{u\sin\theta}$.}
\opt{C}{The magnitude of acceleration of dog at this instant is $\dfrac{uv\cos\theta}{r}$}
\opt{D}{Radius of curvature of path followed by dog at this instant $\dfrac{ur}{v\cos\theta}$}
\end{qoptions}
\qn{8} A small bead of mass m can slide on a frictionless fixed ring of radius r. With the help of two identical strings of force constant k, the bead is connected to two nails A and B each on a diameter at a distance 0.5 r from centre O of the ring. Relaxed length of each string is negligible as compared to the radius of the ring. The bead is given a small velocity at point C. What can you predict about subsequent motion of the bead before any of the string strikes a nail? [ignore gravity]\par
\qfig{p2_phy_s2_q08_1}
\begin{qoptions}
\opt{A}{It will keep moving with its initial speed.}
\opt{B}{It will oscillate simple harmonically about the point C.}
\opt{C}{Its angular momentum about centre O of the ring is conserved.}
\opt{D}{Total elastic potential energy stored in both the strings is 1.25kr$^2$.}
\end{qoptions}
\qn{9} A solid paraboloid of base radius R$_0$ and height h$_0$ is having uniform volume mass density of $\rho_\mathrm s$ is Inverted and placed just above the surface as shown, of a liquid having volume mass density of $\rho_\mathrm L$. When paraboloid is released from rest it has been observed that when it becomes completely submerged in the liquid the paraboloid comes to the rest again. There exist uniform gravity of g. Assume that the liquid body is very large so that liquid level is not changing as the paraboloid enters into the liquid.\par
\qfig{p2_phy_s2_q09_1}
\begin{qoptions}
\opt{A}{The phenomenon is true for all paraboloids irrespective of their dimensions as long as density ratio is maintained.}
\opt{B}{The phenomenon is not true for all paraboloids irrespective of their dimensions as long as density ratio is maintained.}
\opt{C}{The ratio of $\dfrac{\rho_L}{\rho_s}$ is 2.}
\opt{D}{The ratio of $\dfrac{\rho_L}{\rho_s}$ is 3.}
\end{qoptions}
\qn{10} A ideal gas having f degrees of freedom is kept inside a thermally insulated vessel, the vessel has two large pistons each with charge of +Q and $-$Q. Assume charges are uniformly distributed. Pistons can move without friction inside the vessel. At an instant the charge of pistons is increased k times with the help of an external agent. Then which of the following statements are correct\par
\qfig{p2_phy_s2_q10_1}
\begin{qoptions}
\opt{A}{The Ratio of final to initial temperature of gas in equilibrium is $\dfrac{K^2+f}{2+f}$.}
\opt{B}{The Ratio of final to initial temperature of gas in equilibrium is $\dfrac{2K^2+f}{2+f}$.}
\opt{C}{Work done by the agent in the process is 2 (K$^2$ $-$ 1)nRT$_1$ where T$_1$ is initial temperature.}
\opt{D}{Work done by the agent in the process is $\left(K^2-1\right)\dfrac{Q^2L_1}{A\varepsilon_0}$ where L$_1$ is initial separation between pistons.}
\end{qoptions}
\sect{SECTION-II(i)}
\paraq{1}{2}
There is a hollow prism with hexagonal base of side L$_0$ and height H$_0$ as shown. There are identical non conducting uniformly charged solid spheres are placed at each corner and two identical spheres at the centre of flat hexagonal faces. The charges spheres have uniform density of d$_0$ and radius of R$_0$. What will be the flux passing through the shaded region and average charge density ($\phi$) hexagonal prism. It is given that the flux passing through the hexagonal flat surface is k times the total flux passing through the prism.\par
(Take $\dfrac{R_0{}^3d_0}{\varepsilon_0}$ = 1 SI units \& $\dfrac{d_0R_0{}^3}{L_0^2H_0}$ = 1 SI units.)\par
\qfig{p2_phy_s3_para1_2_1}
\qn{1} The value of $\phi_\mathrm s$ is\_\_\_\_\_\_\_
\qn{2} The value of d is\_\_\_\_\_\_\_
\paraq{3}{4}
One side of a thin metal plate is illuminated by the sun. When the air temperature is T$_0$, the temperature of the illuminated side is T$_1$ and that of the opposite side is T$_2$. The temperature of each sides be T$_1'$ \& T$_2'$, if another plate of double thickness is used. (Take T$_0$ = 300 K, T$_1$ = 330 K, T$_2$ = 320 K )
\qn{3} The value of T$_1'$ (K) is \_\_\_\_\_\_\_\_\_\_\_
\qn{4} The value of T$_2'$ (K) is \_\_\_\_\_\_\_\_\_
\paraq{5}{6}
A boat is travelling in a river with a speed of 10 m/s along the stream flowing with a speed of 2 m/s. From this boat, a sound transmitter is lowered into the river through a rigid support. The wavelength of the sound emitted from the transmitter inside the water is 14.45 mm. Assume that attenuation of sound in water and air is negligible. The frequency detected by a receiver kept inside the river downstream is f$_0$ and the transmitter and the receiver are now pulled up into air. The air is blowing with a speed 5 m/sec in the direction opposite the river stream the frequency of the sound detected by the receiver is f$_1$. (Temperature of the air and water = 20$^\circ$C; Density of river water = 10$^3$ kg/m$^3$; Bulk modulus of the water = 2.088 $\times$10$^9$ Pa; Gas constant R = 8.31J/mol $-$K ; Mean molecular mass of air = 28.8 $\times$ 10$^{-3}$ kg/mol; Cp/Cv for air = 1.4)\par
Note: Boat velocity is with respect to ground \& receiver is stationary w.r.t ground
\qn{5} The value of f$_0$ (KHZ) is\_\_\_\_\_\_\_\_\_\_
\qn{6} The value of f$_1$ (KHZ) is\_\_\_\_\_\_\_\_\_\_
\sect{SECTION-II(ii)}
\qn{7} A conducting sphere with inner radius R and outer radius 3R has some alloys mixed in it due to which its resistivity changes with radial distance r according to the function $\rho_r$ = Kr, where K is a +ve constant. Inner surface of the sphere is grounded while outer surface of sphere in maintained at potential V$_0$ with the help of a battery of emf V$_0$. Battery is ideal. The potential of the sphere at radial distance of 2R is V$_0\dfrac{\ln a}{\ln b}$. Find (a + b).\par
\qfig{p2_phy_s4_q07_1}
\qn{8} In the circuit shown in the figure the electromotive force of the battery is 9 V and its internal resistance is 15 $\Omega$.The two identical voltmeters can be considered ideal. Let V$_1$ and V$'_1$ reading of 1$^\mathrm{st}$ voltmeter when switch is open and closed respectively. Similarly, V$_2$ and V$'_2$ be the reading of 2$^\mathrm{nd}$ voltmeter when switch is open and closed respectively. Then find the value of $\dfrac{V_2'-V_2}{V_1-V_1'}$.\par
\qfig{p2_phy_s4_q08_1}
\qn{9} The given thin lenses have R$_1$ = R$_2$ = 10 cm and $\mu_1$ = 1, $\mu_2$ = 2 and $\mu_3$ = 4. If final image formed at object and value of `d' is x $\times$ 10 cm . Find `x'.\par
\qfig{p2_phy_s4_q09_1}
\partsec{2}{CHEMISTRY}
\sect{SECTION-I(i)}
\qn{1}\par
\qfig{p2_chem_s1_q01_1}
The compound (B) would be
\begin{qoptions}
\opt{A}{\qfigi{p2_chem_s1_q01_A}}
\opt{B}{\qfigi{p2_chem_s1_q01_B}}
\opt{C}{\qfigi{p2_chem_s1_q01_C}}
\opt{D}{\qfigi{p2_chem_s1_q01_D}}
\end{qoptions}
\qn{2} How many methoxy group present in the molecule of a compound (A) would be present by the result of a Zeisel method. Treatment of compound (A) C$_{20}$H$_{11}$O$_4$N with hot conc. HI gives CH$_3$I indicating presence of methoxy group. When 4.24 mg of a compound (A) is treated with hot conc. HI and CH$_3$I thus produced is digested with alc. AgNO$_3$, 11.62 mg of AgI is obtained.
\begin{qoptions}
\opt{A}{1}
\opt{B}{2}
\opt{C}{3}
\opt{D}{4}
\end{qoptions}
\qn{3} The blue colour produced on adding H$_2$O$_2$ to acidified K$_2$Cr$_2$O$_7$ solution and shaking with ether is due to the formation of
\begin{qoptions}
\opt{A}{Cr$_2$O$_3$}
\opt{B}{CrO$_4^{-2}$}
\opt{C}{CrO$_3$}
\opt{D}{CrO$_5$}
\end{qoptions}
\qn{4} A 0.001 molal solution of a complex [MA$_8$] in water has the freezing point $-$0.0054$^\mathrm o$C . Assuming 100\% ionization of the complex salt and K$_\mathrm f$ for H$_2$O = 1.86 K-Kgmol$^{-1}$, write the correct representation of the complex.
\begin{qoptions}
\opt{A}{[MA$_8$]}
\opt{B}{[MA$_7$]A}
\opt{C}{[MA$_6$]A$_2$}
\opt{D}{[MA$_5$]A$_3$}
\end{qoptions}
\sect{SECTION-I(ii)}
\qn{5} Cation(s) which form an ammine complex ion with excess of NH$_3$ is/are
\begin{qoptions}
\opt{A}{Ag$^+$}
\opt{B}{Cu$^{+2}$}
\opt{C}{Co$^{+2}$}
\opt{D}{Al$^{+3}$}
\end{qoptions}
\qn{6} Which of the following statements regarding P$_4$O$_{10}$ is correct ?
\begin{qoptions}
\opt{A}{It is an anhydride of metaphosphoric acid}
\opt{B}{It is a strong dehydrating agent}
\opt{C}{In P$_2$O$_5$, each phosphorus is attached to three oxygen atoms through single bond and also with another oxygen atom through double bond.}
\opt{D}{The dehydrated product of HNO$_3$ obtained by treating it with P$_4$O$_{10}$ is NO$_2$.}
\end{qoptions}
\qn{7} Choose the correct option(s) for the following set of reactions\par
\qfig{p2_chem_s2_q07_1}
\begin{qoptions}
\opt{A}{\qfigi{p2_chem_s2_q07_A}}
\opt{B}{\qfigi{p2_chem_s2_q07_B}}
\opt{C}{\qfigi{p2_chem_s2_q07_C}}
\opt{D}{\qfigi{p2_chem_s2_q07_D}}
\end{qoptions}
\qn{8} The reaction(s) leading to the formation of 1, 3, 5 trimethyl benzene is/are
\begin{qoptions}
\opt{A}{\qfigi{p2_chem_s2_q08_A}}
\opt{B}{\qfigi{p2_chem_s2_q08_B}}
\opt{C}{\qfigi{p2_chem_s2_q08_C}}
\opt{D}{\qfigi{p2_chem_s2_q08_D}}
\end{qoptions}
\qn{9} Which of the following statement(s) is/are correct?
\begin{qoptions}
\opt{A}{If temperature coefficient is greater than zero, cell reaction is endothermic.}
\opt{B}{If temperature coefficient is less than zero, cell reaction is endothermic.}
\opt{C}{If temperature coefficient is less than zero, cell reaction is exothermic.}
\opt{D}{If E$_\mathrm{cell}$ is negative then $\Delta$G is negative and cell reaction is spontaneous.}
\end{qoptions}
\qn{10} In which of the following pair of solution will the values of Vant Hoff factor be same? (Assuming same dissociation)
\begin{qoptions}
\opt{A}{0.05 M K$_4$ [Fe(CN)$_6$] and 0.10 M FeSO$_4$}
\opt{B}{0.10 M K$_4$ [Fe(CN)$_6$] and 0.05 M FeSO$_4$.(NH$_4$)$_2$ SO$_4$.6H$_2$O}
\opt{C}{0.20 M NaCl and 0.10 M BaCl$_2$}
\opt{D}{0.05 M FeSO$_4$.(NH$_4$)$_2$ SO$_4$.6H$_2$O and 0.02 M KCl.MgCl$_2$.6H$_2$O}
\end{qoptions}
\qfullbegin
\sect{SECTION-II(i)}
\paraq{1}{2}
For the following reaction scheme\par
\qfig{p2_chem_s3_para1_2_1}
\qfullend
\qn{1} The molecular weight of B is \dots\dots\dots\dots\dots
\qn{2} The molecular weight of E is\dots\dots\dots\dots.
\paraq{3}{4}
A 10$^{-3}$ molal aqueous solution of a chromium complex having same number of ammonia molecule and Cl ions freezes at $-$0.00558$^\circ$C. If all the ammonia molecules are acting as ligands and the percent of hydrogen in the complex is 4.6 (Given K$_\mathrm F$(H$_2$O) = 1.86K Kgmol$^{-1}$ ) and assume 100\% ionization self.
\qn{3} Find the number of Cl atom is complex is\_\_\_.
\qn{4} The value of Vant Hoff factor is\_\_\_\_\_\_.
\paraq{5}{6}
The isoelectric point (P$_\mathrm I$) is the P$_\mathrm H$ at which the amino acid exists only as a dipolar ion with net charge zero. Structure of lysine and aspartic acid are $\mathrm{C\overset{\ominus}{O}O\,CH\left(\overset{+}{N}H_3\right)(CH_2)_4NH_2}$ and $\mathrm{CO\overset{\ominus}{O}\,CH\left(\overset{+}{N}H_3\right)CH_2COOH}$ respectively. The P$_{\mathrm{Ka}_1}$, P$_{\mathrm{Ka}_2}$ and P$_{\mathrm{Ka}_3}$ of the dictation lysine are 2.18, 8.95 and 10.53 respectively. The P$_{\mathrm{Ka}_1}$, P$_{\mathrm{Ka}_2}$ and P$_{\mathrm{Ka}_3}$ of the cation of aspartic acid are 1.88, 3.65 and 9.60 respectively.
\qn{5} Isoelectric point lysine is\_\_\_\_\_\_\_\_\_\_.
\qn{6} Isoelectric point of Aspartic acid is \_\_\_\_\_\_\_\_\_\_.
\sect{SECTION-II(ii)}
\qn{7} Find the quantum numbers of the excited state of electrons (Let n$_1$ and n$_2$) in He$^+$ ion which on transition to ground state emit two photons of wavelength 30.4 nm and 108.5 nm respectively. Find the sum of (n$_1$ + n$_2$) = ?
\qn{8} A compound AB forms NaCl type crystals and another compound XY forms CsCl type crystals. The formula. Mass of XY is three times that of AB but the cubic edge length of unit cell of AB crystals is twice that of XY. Calculate the ratio of density of XY to that of AB.
\qn{9} Size of ring in product formed C is ?\par
\qfig{p2_chem_s4_q09_1}
\partsec{3}{MATHEMATICS}
\sect{SECTION-I(i)}
\qn{1} Polynomial function f (x) satisfying the condition $f(x)f\left(\dfrac1x\right)=f(x)+f\left(\dfrac1x\right)$. If f (10) = 1001, then f(20) is
\begin{qoptions}
\opt{A}{7001}
\opt{B}{8001}
\opt{C}{8000}
\opt{D}{None of these}
\end{qoptions}
\qn{2} Let $f(x)=e^{\left\{e^{|x|}\,\mathrm{sgn}\,x\right\}}$ and $g(x)=e^{\left[e^{|x|}\,\mathrm{sgn}\,x\right]}$, x $\in$ R where \{ \} and [ ] denotes the fractional and integral part functions, respectively. Also h(x) = log(f(x)) + log (g(x)) then for real x, h(x) is
\begin{qoptions}
\opt{A}{An odd function}
\opt{B}{An even function}
\opt{C}{Neither odd nor an even function}
\opt{D}{Both odd as well as even function}
\end{qoptions}
\qn{3} The length and foot of the perpendicular from the point (2, $-$1, 5) to the line $\dfrac{x-11}{10}=\dfrac{y+2}{-4}=\dfrac{z+8}{-11}$ are
\begin{qoptions}
\opt{A}{$\sqrt{14}$, (1, 2, $-$3)}
\opt{B}{$\sqrt{14}$, (1, $-$2, 3)}
\opt{C}{$\sqrt{14}$, (1, 2, 3)}
\opt{D}{None of these}
\end{qoptions}
\qn{4} If x = 9 is the chord of contact of the hyperbola x$^2$ $-$ y$^2$ = 9, then the equation of the corresponding pair of tangent is
\begin{qoptions}
\opt{A}{9x$^2$ $-$ 8y$^2$ + 18x $-$ 9 = 0}
\opt{B}{9x$^2$ $-$ 8y$^2$ $-$ 18x + 9 = 0}
\opt{C}{9x$^2$ $-$ 8y$^2$ $-$ 18x $-$ 9 = 0}
\opt{D}{9x$^2$ $-$ 8y$^2$ + 18x + 9 = 0}
\end{qoptions}
\sect{SECTION-I(ii)}
\qn{5} Which of the following options are correct ?
\begin{qoptions}
\opt{A}{\resizebox{\linewidth}{!}{$\left[\left({}^n\mathrm C_0+{}^n\mathrm C_3+{}^n\mathrm C_6+\dots..\right)-\dfrac12\left({}^n\mathrm C_1+{}^n\mathrm C_2+{}^n\mathrm C_4+{}^n\mathrm C_5+\dots..\right)\right]^2$}\par $+\dfrac34\left({}^n\mathrm C_1-{}^n\mathrm C_2+{}^n\mathrm C_4-{}^n\mathrm C_5+\dots.\right)^2=1$}
\opt{B}{If a and b are two positive numbers such that a$^5$b$^2$ = 4 then the maximum value of $\log_{2^{1/5}}\left(a^2\right).\log_{2^{1/2}}\left(b^2\right)$ is equal to 4}
\opt{C}{Constant team in $((((((x-2)^2-2)^2-2)^2-2)^2.......2)^2$ is equal to 2}
\opt{D}{The coefficient of x$^{24}$ in\par\resizebox{\linewidth}{!}{$\left(\dfrac{{}^{25}\mathrm C_1}{{}^{25}\mathrm C_0}-x\right)\left(x-2^2\dfrac{{}^{25}\mathrm C_2}{{}^{25}\mathrm C_1}\right)\left(x-3^2\dfrac{{}^{25}\mathrm C_3}{{}^{25}\mathrm C_2}\right)\times\left(x-4^2\dfrac{{}^{25}\mathrm C_4}{{}^{25}\mathrm C_3}\right)\dots\dots$}\par
\resizebox{0.3\linewidth}{!}{$\left(x-25^2\dfrac{{}^{25}\mathrm C_{25}}{{}^{25}\mathrm C_{24}}\right)$} is equal to 2925}
\end{qoptions}
\qn{6} a$_1$, a$_1$, a$_2$,...... are distinct terms of an A.P. We call (p,q,r) an increasing triad if a$_\mathrm p$, a$_\mathrm q$, a$_\mathrm r$ are in G.P. where p,q,r $\in$ N such that p $<$ q $<$ r. If (5, 9, 16) is an increasing triad, then which of the following options is/are correct
\begin{qoptions}
\opt{A}{If a$_1$ is a multiple of 4 then every term of the A.P. is an integer}
\opt{B}{(85,149,261) is an increasing triad}
\opt{C}{If the common difference of the A.P. is $\dfrac14$, then its first term is $\dfrac13$}
\opt{D}{ratio of (4k + 1)th term and 4k$^\mathrm{th}$ term can be 4}
\end{qoptions}
\qn{7} If $\displaystyle\int\dfrac{e^{x-1}}{\left(x^2-5x+4\right)}2x\,dx=AF\,(x-1)+BF\,(x-4)+C$ and $F(x)=\displaystyle\int\dfrac{e^x}{x}dx$, then
\begin{qoptions}
\opt{A}{A = $-$ 2/3}
\opt{B}{B =(4/3)e$^3$}
\opt{C}{A = 2 / 3}
\opt{D}{B = (8 / 3)e$^3$}
\end{qoptions}
\qn{8} Let OA, OB, OC be coterminous edges of a cuboid. If l, m, n be the shortest distances between the sides OA, OB, OC and their respective skew body diagonals to them respectively, then the value of $\dfrac{\left(\frac1{l^2}+\frac1{m^2}+\frac1{n^2}\right)}{\left(\frac1{OA^2}+\frac1{OB^2}+\frac1{OC^2}\right)}$ is less than
\begin{qoptions}
\opt{A}{2}
\opt{B}{3}
\opt{C}{4}
\opt{D}{5}
\end{qoptions}
\qn{9} Let the equation of a straight line L in complex form be $a\bar z+\bar az+b=0$, where a is a complex number and b is a real number, then
\begin{qoptions}
\opt{A}{the straight line $\dfrac{z-c}a+\dfrac{i(\bar z-\bar c)}a=0$ makes an angle of 45$^\circ$ with L and passes through a point c (where c is a complex number)}
\opt{B}{the straight line $\dfrac{z-c}a=\dfrac{i(\bar z-\bar c)}a=0$ makes an angle of 45$^\circ$ with L and passes through a point c (where c is a complex number)}
\opt{C}{the complex slope of the line L is $-\dfrac aa$}
\opt{D}{the complex slope of the line L is $\dfrac aa$}
\end{qoptions}
\qn{10} The function f(x) = x$^{1/3}$ (x$-$1)
\begin{qoptions}
\opt{A}{has two inflection points}
\opt{B}{has one point of extremum}
\opt{C}{is non $-$ differentiable}
\opt{D}{Range of f(x) is [$-$ 3 $\times$ 2$^{-8/3}$, $\infty$)}
\end{qoptions}
\sect{SECTION-II(i)}
\paraq{1}{2}
\textbf{Question Stem}\par
Consider the locus of the complex number z in the Argand plane given by Re(z) $-$ 2 = $|z-7+2i|$. Let P($z_1$), Q($z_2$). Be two complex numbers satisfying the given locus and also satisfying $\arg\left(\dfrac{z_1-(2+\alpha i)}{z_2-(2+\alpha i)}\right)=\dfrac\pi2$ ($\alpha\in$ R).
\qn{1} The minimum value of PR.QR where R represents the point (7, $-$2) is
\qn{2} $\arg\left(\dfrac{z_1-(7-2i)}{z_1-(7-2i)}\right)$ equals
\paraq{3}{4}
\textbf{Question Stem}\par
T is the region of the plane x + y + z = 1 with x, y, z $>$0. S is the set of points (a, b, c) in T such that just two of the following three inequalities hold : $a\leqslant\dfrac12$, $b\leqslant\dfrac13$, $c\leqslant\dfrac16$.
\qn{3} Area of the region T is
\qn{4} Area of the region S is
\paraq{5}{6}
\textbf{Question Stem}\par
Consider the line $\dfrac{x-1}2=\dfrac y{-1}=\dfrac{z+1}2$ and the point C($-$1, 1, 2). Let D be the image of C in the line.
\qn{5} The distance C from the line is
\qn{6} The distance of the origin of the plane through the line and the point C is
\sect{SECTION-II(ii)}
\qn{7} A, B, C and D cut a pack of 52 cards successively in the order given. If the person who cuts a spade first receives Rs. 350 and if the expectation of A is $\lambda$ then $\left[\dfrac\lambda{64}\right]$ is equal to (where [.] denotes greatest integer function)\_\_\_\_\_\_\_\_\_
\qn{8} Let f(x) = $\displaystyle\lim_{x\to\infty}\dfrac{x^{2n-1}+ax^2+bx}{x^{2n}+1}$. If f(x) is continuous for all x $\in$ R, then the value of a + 8b is
\qn{9} A and B are two independent events. The probability that both A and B occur is 1/6 and the probability that at least one of them occurs is 2/3. Find 8P(A) + 9P(B) if P(A) $>$ P(B).
\end{multicols}
\end{document}
